\documentclass[12pt,a4paper]{article}
\usepackage[utf8]{inputenc}
\usepackage[T1]{fontenc}
\usepackage[brazil]{babel}
\usepackage{lmodern}
\usepackage{geometry}
\geometry{top=2.5cm,bottom=2.5cm,left=3cm,right=2cm}
\usepackage{setspace}
\onehalfspacing
\usepackage{graphicx}
\usepackage{booktabs}
\usepackage{amsmath,amssymb}
\usepackage{float}
\usepackage{enumitem}
\usepackage{hyperref}
\usepackage{xcolor}
\usepackage{listings}
\usepackage[alf]{abntex2cite}

\hypersetup{colorlinks=true,linkcolor=black,citecolor=black,urlcolor=blue}
\lstset{basicstyle=\ttfamily\small,breaklines=true,frame=single,numbers=left,
numberstyle=\tiny,showstringspaces=false,captionpos=b}

% Edite os dados do projeto
\newcommand{\TituloProjeto}{Título do projeto}
\newcommand{\Disciplina}{Resolução de Problemas via Modelagem Matemática}
\newcommand{\Universidade}{Universidade Federal de São Paulo (UNIFESP)}
\newcommand{\Ano}{2026}

\begin{document}
\begin{titlepage}
    \centering
    {\large \Universidade\par}
    \vspace{1.5cm}
    {\large \Disciplina\par}
    \vspace{2.5cm}
    {\LARGE\bfseries \TituloProjeto\par}
    \vspace{1.5cm}
    {\large Relatório de projeto\par}
    \vfill
    \begin{minipage}{0.8\textwidth}
        \centering
        \textbf{Integrantes do grupo (ordem alfabética)}\\[0.4cm]
        Nome Sobrenome\\
        Nome Sobrenome\\
        Nome Sobrenome
    \end{minipage}
    \vfill
    {\large São José dos Campos\par}
    {\large \Ano\par}
\end{titlepage}

\tableofcontents
\newpage

\section{Projeto}
Apresente o contexto, a relevância e o escopo do projeto. Informe o título e os integrantes do grupo em ordem alfabética.

\subsection{Proponente}
Identifique a empresa, instituição ou pessoa que interage com o grupo. Descreva brevemente sua atuação e sua relação com o problema.

\subsection{Definição do problema e objetivo}
Descreva o problema, o contexto, as decisões envolvidas, as restrições relevantes e o objetivo do trabalho. Delimite o escopo e os critérios de avaliação, quando aplicável.

\subsection{Fundamentos teóricos}
Apresente os conceitos, teorias, aplicações, algoritmos e técnicas de modelagem pertinentes. Explique a abordagem adotada e cite as fontes ao longo do texto.

\section{Modelo matemático}
Apresente o modelo proposto e explique sua relação com os fundamentos teóricos. Defina conjuntos, índices, parâmetros, variáveis de decisão, função objetivo e restrições, conforme aplicável.

\subsection{Conjuntos e índices}
\begin{itemize}[leftmargin=*]
    \item $\mathcal{I}$: descrição do conjunto;
    \item $i \in \mathcal{I}$: interpretação do índice.
\end{itemize}

\subsection{Parâmetros de entrada}
\begin{table}[H]
\centering
\caption{Parâmetros do modelo}
\begin{tabular}{lll}
\toprule
Símbolo & Descrição & Unidade/fonte\\
\midrule
$p_1$ & Descrição do parâmetro & Unidade ou fonte\\
$p_2$ & Descrição do parâmetro & Unidade ou fonte\\
\bottomrule
\end{tabular}
\end{table}

\subsection{Variáveis de decisão}
\begin{table}[H]
\centering
\caption{Variáveis de decisão}
\begin{tabular}{lll}
\toprule
Símbolo & Descrição & Domínio\\
\midrule
$x_1$ & Descrição da variável & $\{0,1\}$, $\mathbb{R}_+$ etc.\\
$x_2$ & Descrição da variável & $\mathbb{Z}_+$ etc.\\
\bottomrule
\end{tabular}
\end{table}

\subsection{Formulação}
Escreva a função objetivo e as restrições. Explique a interpretação de cada expressão e as hipóteses adotadas.
\begin{align}
\min \quad & f(x) \label{eq:objetivo}\\
\text{sujeito a} \quad & g_j(x) \leq 0, \quad j \in \mathcal{J}, \label{eq:restricao}\\
& x \in X. \nonumber
\end{align}

\section{Resultados}
Apresente resultados analíticos ou computacionais. Inclua tabelas, gráficos, imagens, listas e medidas de desempenho, quando pertinentes. Descreva os dados e as condições experimentais. Discuta limitações.

\subsection{Resultados quantitativos}
\begin{table}[H]
\centering
\caption{Exemplo de tabela de resultados}
\begin{tabular}{lrr}
\toprule
Instância & Valor objetivo & Tempo (s)\\
\midrule
Instância 1 & 0,00 & 0,00\\
Instância 2 & 0,00 & 0,00\\
\bottomrule
\end{tabular}
\end{table}

\subsection{Gráficos e imagens}
Envie a figura ao projeto Overleaf e atualize o nome do arquivo.
\begin{figure}[H]
\centering
% Exemplo: envie figura.png ao Overleaf.
\includegraphics[width=0.75\textwidth]{figura.png}
\caption{Descrição da figura.}
\label{fig:exemplo}
\small Fonte: elaboração própria.
\end{figure}

\subsection{Discussão e limitações}
Interprete os resultados em relação aos objetivos. Discuta limitações do modelo, dos dados, dos algoritmos ou da implementação.

\section{Oportunidades para pesquisas futuras}
Indique extensões, melhorias e novas questões que possam ser investigadas.

\section{Conclusões}
Retome o problema e os objetivos. Sintetize as principais contribuições, os resultados e as limitações mais importantes.

\section{Referências}
Inclua todas as obras citadas no texto. Cadastre as referências no arquivo \texttt{referencias.bib} e cite-as com comandos como \texttt{\textbackslash cite\{chave\}}.
\bibliography{referencias}

\appendix
\section{Documentação complementar}
Inclua documentação necessária para compreender e reproduzir o projeto:
\begin{itemize}[leftmargin=*]
    \item códigos-fonte e modelos computacionais;
    \item fontes web e pessoas consultadas, quando pertinente;
    \item problemas-teste, dados e exemplos executados;
    \item instruções de instalação e execução.
\end{itemize}

\section{Repositório e reprodutibilidade}
Os códigos-fonte, modelos e arquivos necessários devem ser disponibilizados em uma conta de equipe no GitHub. Informe o link do repositório e mantenha um arquivo \texttt{README.md} com instruções de instalação, execução, dados e organização dos arquivos.

\noindent\textbf{Repositório GitHub:} \url{https://github.com/organizacao/repositorio}

\section{Outros materiais}
Inclua outros apêndices ou materiais relevantes.

\end{document}

